%!TEX root = ../csuthesis_main.tex
% 设置中文摘要
\keywordscn{物理信息神经网络\quad 偏微分方程\quad 自适应采样\quad 损失权重\quad 网络结构}
%\categorycn{TP391}
\begin{abstractzh}

物理信息神经网络（Physics-Informed Neural Network, PINN）是一种将深度学习与先验物理知识相结合的机器学习框架，用于解决复杂的物理问题，特别是偏微分方程（Partial Differential Equations, PDE）的求解和系统动力学模拟。然而，在实际应用中，PINN模型的训练过程往往面临着收敛速度慢、计算效率低的问题。

本文首先全面介绍了PINN基础理论及其在处理复杂PDE中的应用，对比了PINN与传统机器学习的区别，并分析了PINN相对于传统数值计算方法的优势。然后针对Navier-Stokes方程描述的漩涡脱落现象进行了具体求解，并从自适应采样、动态调整损失权重、改进网络结构三个方向对PINN进行优化。

在采样方法上，首先介绍了传统的固定采样与随机重采样，然后提出两类自适应采样策略：基于残差的自适应分布方法通过模拟残差较大的区域有效提升了模型预测精度，保留-重采样-释放方法进一步减少了计算负担。在动态调整损失权重上，将神经网络的自适应性和基于残差的注意力融入损失函数，加速了模型收敛进程。在网络结构上，基于Transformer的改进多层感知机方法通过结合残差连接与自注意力机制，提升了模型对复杂数据特征的捕获能力。通过三类方法的融合最终实现了对PINN模型的优化加速。

最后，总结了本文的主要研究结论并针对 PINN 算法的改进提出了若干未来值得深入研究的方向。

\end{abstractzh}